\chapter{Cafeteria Sensorial: um Especialista Conversacional em Preparo de Cafés}
\begin{flushright}
Marcílio Pereira Aguiar
\end{flushright}

A Cafeteria Sensorial é um sistema especialista conversacional para orientação de preparo e interpretação de avaliações sensoriais de cafés especiais. Seu objetivo é apoiar usuários que desejam relacionar características dos grãos, parâmetros de preparo e resultados percebidos. A solução evolui um projeto de serviços web desenvolvido anteriormente, acrescentando conhecimento de domínio, cache semântico e avaliação automatizada.

O assistente não substitui uma avaliação profissional nem certifica qualidade. Recomendações são hipóteses fundamentadas nas informações disponíveis: uma percepção de acidez, por exemplo, não determina isoladamente a causa de um resultado desagradável. O recorte combina orientações pedagógicas de preparo e análise sensorial com consulta a registros locais, distinguindo descritores, preferências e evidências \cite{cafe_sca_ratio,cafe_compass,cafe_sca_cva}.

\section{Arquitetura da Cafeteria Sensorial}

A aplicação mantém três serviços REST independentes, responsáveis por grãos, métodos e degustações. Cada serviço possui contêiner e banco próprios em uma instância PostgreSQL compartilhada. O cliente web consulta os serviços com tempos limites e tentativas independentes; uma API indisponível não interrompe os demais painéis. O servidor MCP funciona como interface entre as consultas REST e o núcleo conversacional.

O núcleo utiliza um LLM Gemini com instruções de domínio e uma base JSON de orientações e fontes. O histórico recente constitui o contexto de trabalho. O modelo decide quando obter informações externas e recebe os resultados das ferramentas antes de elaborar a resposta. Para registros do balcão, as instruções exigem consulta, identificação das evidências e explicitação de indisponibilidades. Para problemas incompletos, o assistente deve solicitar informações e propor testes controlados.

\begin{figure}[htbp]
\centering
\fbox{\parbox{0.91\textwidth}{\centering
Usuário $\rightarrow$ Cliente web $\rightarrow$ Núcleo especialista\\[3pt]
Núcleo: contexto + conhecimento + cache semântico + LLM\\[3pt]
LLM $\leftrightarrow$ MCP $\rightarrow$ APIs de grãos, métodos e degustações\\[3pt]
APIs $\rightarrow$ bancos separados no PostgreSQL\\[3pt]
Script de avaliação $\rightarrow$ chat completo $\rightarrow$ LLM julgador
}}
\caption{Componentes e fluxo de consulta da Cafeteria Sensorial.}
\label{fig:cafe-arquitetura}
\end{figure}

O cache utiliza \textit{SentenceTransformers} para representar perguntas e ChromaDB para persistir vetores e respostas. Antes de gerar uma nova resposta, procura-se uma pergunta próxima por distância de cosseno. O limiar inicial de 0,20 acompanha o exemplo da disciplina e deve ser calibrado, não interpretado como probabilidade de acerto. Apenas perguntas conceituais independentes e respostas sem ferramentas são elegíveis. Histórico, números, negações e referências a dados mutáveis impedem a reutilização conservadoramente. Modelo, versão do conhecimento e prazo de validade também limitam o reaproveitamento.

\section{Ferramentas MCP}

O protocolo MCP padroniza a descoberta e invocação das capacidades externas \cite{cafe_mcp}. A solução conserva seis ferramentas, apresentadas na Tabela~\ref{tab:cafe-tools}. Elas consultam as APIs, sem acesso direto aos bancos, e retornam falhas por domínio. O cliente MCP e o servidor utilizam o SDK oficial e transporte HTTP.

\begin{table}[htbp]
\centering
\begin{tabular}{ll}
\hline
Domínio & Ferramentas de consulta\\
\hline
Grãos & \texttt{graos\_listar}, \texttt{graos\_obter}\\
Métodos & \texttt{metodos\_listar}, \texttt{metodos\_obter}\\
Degustações & \texttt{degustacoes\_listar}, \texttt{degustacoes\_obter}\\
\hline
\end{tabular}
\caption{Ferramentas disponibilizadas ao especialista.}
\label{tab:cafe-tools}
\end{table}

Consultas permitem comparar avaliações e relacionar identificadores aos catálogos. As notas finais usam escala de zero a dez; acidez, corpo e doçura, de um a cinco. Médias descrevem a amostra cadastrada e devem informar sua quantidade, sem estabelecer superioridade universal. Os dados iniciais são fictícios. Massas de café, água e bebida podem ser registradas opcionalmente; ausência desses valores não comprova a proporção realmente utilizada.

\section{Desenvolvimento e Avaliação}

A implementação utiliza Python, FastAPI, Pydantic, PostgreSQL, JavaScript, Nginx e Docker Compose. O SDK Google GenAI realiza a comunicação com Gemini; a integração MCP executa as ferramentas escolhidas pelo modelo. A interface identifica a origem da resposta e permite iniciar nova conversa. Falha no cache mantém o atendimento pelo LLM, com indicação de que o cache está indisponível.

O script de avaliação aplica \textit{LLM-as-a-Judge} ao endpoint completo do chat, com cache desativado. Doze casos cobrem conceitos, cálculo de proporção, diagnóstico incompleto, catálogo, médias, métodos, dados ausentes, registro inexistente, escopo e tentativa de fabricação. O julgador recebe pergunta, resposta, critérios esperados e evidências MCP. Avalia correção, fundamentação, clareza, tratamento da incerteza e respeito ao escopo, produzindo decisão estruturada e justificativa.

Dois controles com respostas propositalmente incorretas devem ser rejeitados antes dos casos reais. Aprovação exige veredicto exato e coerência entre todos os critérios; ferramentas obrigatórias também são verificadas pelo script. O relatório separa aprovações, reprovações e erros e inclui casos planejados não concluídos no denominador. Modelos e versão do conhecimento são registrados, permitindo identificar as condições da execução.

% BEGIN VALIDATION
Em uma execução real, os dois controles do julgador foram registrados e 12 de 12 casos planejados foram aprovados, com 0 reprovações e 0 erros. O estado da execução foi completed. O ensaio de cache recuperou uma paráfrase com distância de cosseno 0.027, inferior ao limiar 0,20. A integração REST/MCP e a recuperação dos três serviços foram verificadas com os contêineres e o navegador reais. Esses resultados descrevem apenas os casos e as condições registrados, sem estabelecer acurácia geral. Assistente e julgador utilizaram o mesmo alias de modelo com instruções distintas, limitando a independência da avaliação.
% END VALIDATION

Os testes de software verificam contrato, integração, falhas e isolamento do cache; modelos simulados não demonstram qualidade técnica real. O julgamento automático também pode errar e compartilha limitações dos LLMs. A revisão humana das justificativas, a pequena cobertura de casos e a dependência de cotas externas devem acompanhar a interpretação dos resultados.

\section{Reproducibilidade e Reuso}

A inicialização é realizada por \texttt{docker compose up} \texttt{-d --build}. O script \texttt{julgamento.py} é executado no contêiner do chat e grava relatório JSON persistente. A configuração de modelos, cache e credenciais é feita por variáveis de ambiente; chaves não integram o material distribuído.

O \href{https://doi.org/10.5281/zenodo.24386960}{código-fonte e os relatórios de avaliação} estão disponíveis no Zenodo. O \href{https://zenodo.org/records/24387783}{vídeo demonstrativo} apresenta a inicialização do sistema, as consultas via MCP, o uso do cache semântico e a execução do julgador.
